\documentclass[12pt,oneside]{scrbook}


% Fonts and typography

%% Typography
\usepackage[no-math]{fontspec}
\defaultfontfeatures{Scale = 1}

%% Fonts
\setmainfont[Ligatures=TeX]{Verdana}
\setsansfont[Ligatures=TeX]{Cambria Bold Italic}
\setmonofont{Consolas}

% Section fonts
\usepackage{sectsty}
\allsectionsfont{\sffamily}
\subsectionfont{\sffamily}
\chapterfont{\LARGE\sffamily}


%% Set polyglossia language
\usepackage{polyglossia}
\setdefaultlanguage{english}


% Page

%% Use full page in book style
\usepackage{fullpage}

%% Set line spacing
\usepackage{setspace}
\setstretch{1.2}

%% Disable paragraph indentation
\usepackage{parskip}

% Colors

\usepackage{xcolor}

%% Tango color scheme
\definecolor{SkyBlue}{HTML}{3465A4}
\definecolor{DarkSkyBlue}{HTML}{204A87}

\definecolor{Plum}{HTML}{75507B}

\definecolor{ScarletRed}{HTML}{CC0000}


\definecolor{Aluminium1}{HTML}{EEEEEC}
\definecolor{Aluminium6}{HTML}{2e3436}

\definecolor{Black}{HTML}{000000}

% Use fancy chapter numbering and set the color
\usepackage{quotchap}
\definecolor{chaptergrey}{HTML}{204A87}

% Listings
\usepackage{upquote}

\usepackage{listings}

\lstdefinelanguage{JavaScript}{
  keywords = {typeof, new, true, false, catch, function, return, null, catch, switch, var, if, in, while, do, else, case, break},
  keywordstyle = \color{SkyBlue}\bfseries,
  ndkeywords = {class, export, boolean, throw, implements, import, this},
  ndkeywordstyle = \color{Aluminium6}\bfseries,
  identifierstyle = \color{Black},
  sensitive = false,
  comment = [l]{//},
  morecomment = [s]{/*}{*/},
  commentstyle = \color{Plum}\ttfamily,
  stringstyle = \color{ScarletRed}\ttfamily,
  morestring = [b]',
  morestring = [b]"
}

\lstset{
  language = JavaScript,
  backgroundcolor = \color{Aluminium1},
  extendedchars = true,
  basicstyle = \normalsize\ttfamily,
  showstringspaces = false,
  showspaces = false,
  tabsize = 1,
  breaklines = true,
  showtabs = false
}


% Links

%% Hyperref
\usepackage[colorlinks, breaklinks, bookmarks, xetex]{hyperref}

\hypersetup {
  pdfauthor={Leslie J. Newville},
  pdftitle={Development of the Phonograph at Alexander Graham Bell's Volta Laboratory},
  linkcolor = DarkSkyBlue,
  citecolor = DarkSkyBlue,
  filecolor = DarkSkyBlue,
  urlcolor = DarkSkyBlue
}

%% Don’t use Mono font for URLs
\urlstyle{same}


% Images

\usepackage{graphicx}
\usepackage{incgraph}

\usepackage{eso-pic}
\newcommand\Cover{%
    \put(0,0){%
        \parbox[b][\paperheight]{\paperwidth}{%
        \vfill
        \centering
        \includegraphics[height=\paperheight,width=\paperwidth]{cover.jpg}
        \vfill
}}}


% Pandoc hacks

%% Normal enumerates processing
\usepackage{enumerate}

%% Disable section numbers
\setcounter{secnumdepth}{0}

% Makes table of contents look nicer
\usepackage{tocstyle}
\usetocstyle{standard}

% Set level of sections to be shown in Table of Contents
% 0 = chapters, 1 = chapters/sections, 2 = chapters/sections/subsections
\setcounter{tocdepth}{1}

\begin{document}

  % Title page
  \AddToShipoutPicture*{\Cover} % * = only use it on one page
  \ClearShipoutPicture
  
  \null\newpage
  
  \mainmatter
  \phantomsection  
  \addcontentsline{toc}{chapter}{Contents}
  \renewcommand{\contentsname}{Contents}
  \tableofcontents
  
  % Book contents

  \section*{Preface}\label{preface}
  \addcontentsline{toc}{section}{Preface}
  
  The Project Gutenberg eBook, Development of the Phonograph at Alexander
  Graham Bell's Volta Laboratory, by Leslie J. Newville
  
  This eBook is for the use of anyone anywhere at no cost and with almost
  no restrictions whatsoever. You may copy it, give it away or re-use it
  under the terms of the Project Gutenberg License included with this
  eBook or online at www.gutenberg.org
  
  Title: Development of the Phonograph at Alexander Graham Bell's Volta
  Laboratory Contributions from the Museum of History and Technology,
  United States National Museum Bulletin 218, Paper 5, (pages 69-79)
  
  Author: Leslie J. Newville
  
  Release Date: September 27, 2009 {[}eBook \#30112{]}
  
  Language: English
  
  Character set encoding: ISO-646-US (US-ASCII)
  
  \textbf{\emph{START OF THE PROJECT GUTENBERG EBOOK DEVELOPMENT OF THE
  PHONOGRAPH AT ALEXANDER GRAHAM BELL'S VOLTA LABORATORY}}
  
  E-text prepared by Chris Curnow, Joseph Cooper, Stephanie Eason, and the
  Project Gutenberg Online Distributed Proofreading Team
  (http://www.pgdp.net)
  
  Note: Project Gutenberg also has an HTML version of this file which
  includes the original illustrations. See 30112-h.htm or 30112-h.zip:
  (http://www.gutenberg.org/files/30112/30112-h/30112-h.htm) or
  (http://www.gutenberg.org/files/30112/30112-h.zip)
  
  Transcriber's note:
  
  \begin{verbatim}
    Passages in italics are enclosed between underscores
    (_italics_).
  
    Additional spacing after some of the quotations is intentional
    to indicate both the end of a quotation and the beginning of a
    new paragraph as presented in the original text.
  \end{verbatim}
  
  Contributions from the Museum of History and Technology, United States
  National Museum Bulletin 218, Paper 5, (pages 69-79)
  
  DEVELOPMENT OF THE PHONOGRAPH AT ALEXANDER GRAHAM BELL'S VOLTA
  LABORATORY
  
  LESLIE J. NEWVILLE
  
  DEVELOPMENT OF THE PHONOGRAPH AT ALEXANDER GRAHAM BELL'S VOLTA
  LABORATORY
  
  by
  
  LESLIE J. NEWVILLE
  
  \section*{Beginning}\label{beginning}
  \addcontentsline{toc}{section}{Beginning}
  
  \begin{verbatim}
   _The fame of Thomas A. Edison rests most securely on his genius for
   making practical application of the ideas of others. However, it
   was Alexander Graham Bell, long a Smithsonian Regent and friend of
   its third Secretary S. P. Langley, who, with his Volta Laboratory
   associates made practical the phonograph, which has been called
   Edison's most original invention._
  
   THE AUTHOR: _Leslie J. Newville wrote this paper while he was
   attached to the office of the curator of Science and Technology in
   the Smithsonian Institution's United States National Museum._
  \end{verbatim}
  
  The story of Alexander Graham Bell's invention of the telephone has been
  told and retold. How he became involved in the difficult task of making
  practical phonograph records, and succeeded (in association with Charles
  Sumner Tainter and Chichester Bell), is not so well known.
  
  But material collected through the years by the U. S. National Museum of
  the Smithsonian Institution now makes clear how Bell and two associates
  took Edison's tinfoil machine and made it reproduce sound from wax
  instead of tinfoil. They began their work in Washington, D. C., in 1879,
  and continued until granted basic patents in 1886 for recording in wax.
  
  Preserved at the Smithsonian are some 20 pieces of experimental
  apparatus, including a number of complete machines. Their first
  experimental machine was sealed in a box and deposited in the
  Smithsonian archives in 1881. The others were delivered by Alexander
  Graham Bell to the National Museum in two lots in 1915 and 1922. Bell
  was an old man by this time, busy with his aeronautical experiments in
  Nova Scotia.
  
  It was not until 1947, however, that the Museum received the key to the
  experimental ``Graphophones,'' as they were called to differentiate them
  from the Edison machine. In that year Mrs.~Laura F. Tainter donated to
  the Museum 10 bound notebooks, along with Tainter's unpublished
  autobiography.{[}1{]} This material describes in detail the strange
  machines and even stranger experiments which led in 1886 to a greatly
  improved phonograph.
  
  Thomas A. Edison had invented the phonograph in 1877. But the fame
  bestowed on Edison for this startling invention (sometimes called his
  most original) was not due to its efficiency. Recording with the tinfoil
  phonograph is too difficult to be practical. The tinfoil tears easily,
  and even when the stylus is properly adjusted, the reproduction is
  distorted and squeaky, and good for only a few playbacks. Nevertheless
  young Edison, the ``wizard'' as he was called, had hit upon a secret of
  which men had dreamed for centuries.{[}2{]} Immediately after this
  discovery, however, he did not improve it, allegedly because of an
  agreement to spend the next five years developing the New York City
  electric light and power system.
  
  {[}Illustration: Figure 1.--CHARLES SUMNER TAINTER (1854-1940) from a
  photograph taken in San Diego, California, 1919. (\emph{Smithsonian
  photo 42729-A.}){]}
  
  Meanwhile Bell, always a scientist and experimenter at heart, after his
  invention of the telephone in 1876 was looking for new worlds to
  conquer. If we accept Tainter's version of the story, it was through
  Gardiner Green Hubbard that Bell took up the phonograph challenge. Bell
  had married Hubbard's daughter Mabel in 1879. Hubbard was then president
  of the Edison Speaking Phonograph Co., and his organization, which had
  purchased the Edison patent, was having trouble with its finances
  because people did not like to buy a machine which seldom worked well
  and proved difficult for an unskilled person to operate.
  
  In 1879 Hubbard got Bell interested in improving the machine, and it was
  agreed that a laboratory should be set up in Washington. Experiments
  were also to be conducted on the transmission of sound by light, and
  this resulted in the selenium-cell Photophone, patented in 1881. Both
  the Hubbards and the Bells decided to move to the Capital. While Bell
  took his bride to Europe for an extended honeymoon, his associate
  Charles Sumner Tainter, a young instrument maker, was sent on to
  Washington from Cambridge, Massachusetts, to start the
  laboratory.{[}3{]} Bell's cousin, Chichester Bell, who had been teaching
  college chemistry in London, agreed to come as the third associate.
  During his stay in Europe Bell received the 50,000-franc (\$10,000)
  Volta prize, and it was with this money that the Washington project, the
  Volta Laboratory Association,{[}4{]} was financed.
  
  {[}Illustration: Figure 2.--PHOTOGRAPHING SOUND IN 1884. A rare
  photograph taken at Volta Laboratory, Washington, D. C., by J. Harris
  Rogers, a friend of Bell and Tainter (\emph{Smithsonian photo 44312-E}).
  
  A description of the procedure used is found on page 67, of Tainter's
  unpublished autobiography (see footnote 1). There, Tainter quotes
  Chichester Bell as follows:
  
  ``A jet of bichromate of potash solution, vibrated by the voice, was
  directed against a glass plate immediately in front of a slit, on which
  light was concentrated by means of a lens. The jet was so arranged that
  the light on its way to the slit had to pass through the nappe and as
  the thickness of this was constantly changing, the illumination of the
  slit was also varied. By means of a lens \ldots{} an image of this slit
  was thrown upon a rotating gelatine-bromide plate, on which accordingly
  a record of the voice vibrations was obtained.''{]}
  
  Tainter's story, in his autobiography, of the establishment of the
  laboratory, shows its comparative simplicity:
  
  I therefore wound up my business affairs in Cambridge, packed up all of
  my tools and machines, and \ldots{} went to Washington, and after much
  search, rented a vacant house on L Street, between 13th and 14th
  Streets, and fitted it up for our purpose.{[}5{]}\ldots{} The
  Smithsonian Institution sent us over a mail sack of scientific books
  from the library of the Institution, to consult, and primed with all we
  could learn \ldots{} we went to work.{[}6{]}\ldots{} We were like the
  explorers in an entirely unknown land, where one has to select the path
  that seems to be most likely to get you to your destination, with no
  knowledge of what is ahead.
  
  In conducting our work we had first to design an experimental apparatus,
  then hunt about, often in Philadelphia and New York, for the materials
  with which to construct it, which were usually hard to find, and finally
  build the models we needed, ourselves.{[}7{]}
  
  {[}Illustration: Figure 3.--PAGE FROM NOTEBOOK of Charles Sumner
  Tainter, describing an experiment in sound recording. The Tainter
  notebooks, preserved in the U. S. National Museum, describe experiments
  at the Volta Laboratory, in the 1880's. The Graphophone patents of 1886,
  were the result of this research. (\emph{Smithsonian photo 44312}.){]}
  
  The experimental machines built at the Volta Laboratory include both
  disc and cylinder types, and an interesting ``tape'' recorder. The
  records used with the machines and now in the collections of the U. S.
  National Museum, are believed to be the oldest reproducible records
  preserved anywhere in the world. While some are scratched and cracked,
  others are still in good condition.
  
  By 1881 the Volta associates had succeeded in improving an Edison
  tinfoil machine to some extent. Wax was put in the grooves of the heavy
  iron cylinder, and no tinfoil was used. Rather than apply for a patent
  at that time, however, they deposited the machine in a sealed box at the
  Smithsonian, and specified that it was not to be opened without the
  consent of two of the three men. In 1937 Tainter (fig. 1) was the only
  one still living, so the box was opened with his permission.
  
  For the occasion, the heirs of Alexander Graham Bell gathered in
  Washington, but Tainter was too old and too ill to come from San Diego.
  
  The sound vibrations had been indented in the wax which had been applied
  to the Edison phonograph. The following is the text of the recording:
  ``There are more things in heaven and earth, Horatio, than are dreamed
  of in your philosophy. I am a graphophone and my mother was a
  phonograph.'' Remarked Mrs.~Gilbert Grosvenor,{[}8{]} Bell's daughter,
  when the box was opened in 1937, ``That is just the sort of thing father
  would have said. He was always quoting from the classics.''
  
  {[}Figure 4.--PATENT DRAWINGS from U. S. patent 341214, granted May 4,
  1886, to Chichester Bell and C. S. Tainter.{]}
  
  The method of reproduction used on the machine, however, is even more
  interesting than the quotation. Rather than a stylus and diaphragm, a
  jet of air under high pressure was used.
  
  ``This evening about 7 P. M.,'' Tainter noted on July 7, 1881, ``The
  apparatus being ready the valve upon the top of the air cylinder was
  opened slightly until a pressure of about 100 lbs. was indicated by the
  gage. The phonograph cylinder was then rotated, and the sounds produced
  by the escaping air could be heard, and the words understood a distance
  of at least 8 feet from the phonograph.'' The point of the jet is glass,
  and could be directed at a single groove.
  
  {[}Illustration: Figure 5.--EXPERIMENTAL GRAPHOPHONE photographed in
  1884 at the Volta Laboratory. This is similar to one preserved at the
  Smithsonian Institution. (\emph{Smithsonian photo 44312-D.}){]}
  
  The other experimental Graphophones indicate an amazing range of
  experimentation. While the method of cutting a record on wax was the one
  later exploited commercially, everything else seems to have been tried
  at least once.
  
  The following was noted on Wednesday, March 20, 1881: ``A fountain pen
  is attached to a diaphragm so as to be vibrated in a plane parallel to
  the axis of a cylinder--The ink used in this pen to contain iron in a
  finely divided state, and the pen caused to trace a spiral line around
  the cylinder as it turned. The cylinder to be covered with a sheet of
  paper upon which the record is made\ldots{}. This ink \ldots{} can be
  rendered magnetic by means of a permanent magnet. The sounds were to be
  reproduced by simply substituting a magnet for the fountain
  pen\ldots{}.''
  
  The result of these ideas for magnetic reproduction resulted in patent
  341287, granted on May 4, 1886; it deals solely with ``the reproduction,
  through the action of magnetism, of sounds by means of records in solid
  substances.''
  
  {[}Illustration: Figure 6.--ANOTHER EXPERIMENTAL GRAPHOPHONE,
  photographed at the Volta Laboratory in 1884. (\emph{Smithsonian photo
  44312-F.}){]}
  
  The air jet used in reproducing has already been described. Other jets,
  of molten metal, wax, and water, were also tried. On Saturday, May 19,
  1883, Tainter wrote (see fig. 3):
  
  \begin{verbatim}
   Made the following experiment today:
  
   The cylinder of the Edison phonograph was covered with the coating
   of paraffine-wax and then turned off true and smooth.
  
   A cutting style A., secured to the end of a lever B was then
   adjusted over the cylinder, as shown. Lever B was pivoted at the
   points C-D, and the only pressure exerted to force the style into
   the wax was due to the weight of the parts.
  
   Upon the top of A was fixed a small brass disk, and immediately
   over it a sensitive water jet adjusted, so that the stream of water
   at its sensitive part fell upon the center of the brass disk.
  
   The Phonograph cylinder E, was rotated while words and sounds were
   shouted to the support to which the water jet was attached, and a
   record that was quite visible to the unaided eye was the result.
  \end{verbatim}
  
  The tape recorder, an unusual instrument which recorded mechanically on
  a 3/16-inch strip of wax-covered paper, is one of the machines described
  and illustrated in U. S. patent 341214, dated May 4, 1886 (see fig. 4).
  The strip was coated by dipping it in a solution of beeswax and
  paraffine (one part white beeswax, two parts paraffine, by weight), then
  scraping one side clean and allowing the other side to harden.
  
  {[}Illustration: Figure 7.--ORIGINAL PLANS FOR A DISC GRAPHOPHONE
  PATENTED BY SUMNER TAINTER IN 1888, U. S. PATENT 385886.{]}
  
  The machine of sturdy wood and metal construction, is hand powered by
  means of a knob fastened to the fly wheel. From the fly-wheel shaft
  power is transferred by a small friction wheel to a vertical shaft. At
  the bottom of this shaft a V-pulley transfers motion by belts to
  corresponding V-pulleys beneath the horizontal reels.
  
  The wax strip passes from one 8-inch reel around the periphery of a
  pulley (with guide flanges) mounted above the V-pulleys on the main
  vertical shaft, where it comes in contact with the recording or
  reproducing stylus. It is then taken up on the other reel.
  
  The sharp recording stylus, actuated by a vibrating mica diaphragm, cuts
  the wax from the strip. In reproducing, a dull, loosely mounted stylus,
  attached to a rubber diaphragm, carried sounds through an ear tube to
  the listener.
  
  Both recording and reproducing heads, mounted alternately on the same
  two posts, could be adjusted vertically so that several records could be
  cut on the same 3/16-inch strip.
  
  While this machine was never developed commercially, it is an
  interesting ancestor of the modern tape recorder, which it resembles
  somewhat in design. How practical it was or just why it was built we do
  not know. The tape is now brittle, the heavy paper reels warped, and the
  reproducing head missing. Otherwise, with some reconditioning, it could
  be put into working condition.
  
  Most of the disc machines designed by the Volta associates had the disc
  mounted vertically (see figs. 5 and 6). The explanation is that in the
  early experiments, the turntable, with disc, was mounted on the shop
  lathe, along with the recording and reproducing heads. Later, when the
  complete models were built, most of them featured vertical turntables.
  
  {[}Illustration: Figure 8.--ANOTHER PAGE OF THE PLANS SHOWN IN FIGURE 7.
  The experimental Graphophone built from these plans is in the U. S.
  National Museum (\emph{cat. no. 287665}).{]}
  
  An interesting exception has a horizontal 7-inch turntable (see figs. 7
  and 8). This machine, although made in 1886, is a duplicate of one made
  earlier but taken to Europe by Chichester Bell. Tainter was granted U.
  S. patent 385886 for it on July 10, 1888.
  
  The playing arm is rigid, except for a pivoted vertical motion of 90
  degrees to allow removal of the record or a return to starting position.
  While recording or playing, the record not only rotated, but moved
  laterally under the stylus, which thus described a spiral, recording 150
  grooves to the inch.
  
  The Bell and Tainter records, preserved at the Smithsonian, are both of
  the lateral cut and ``hill-and-dale'' types. Edison for many years used
  the ``hill-and-dale'' method with both cylinder and disc records, and
  Emile Berliner is credited with the invention of the lateral cut
  Gramophone record in 1887. The Volta associates, however, had been
  experimenting with both types, as early as 1881, as is shown by the
  following quotation from Tainter:{[}9{]}
  
  \begin{verbatim}
   The record on the electro-type in the Smithsonian package is of the
   other form, where the vibrations are impressed _parallel_ to the
   surface of the recording material, as was done in the old Scott
   Phonautograph of 1857, thus forming a groove of uniform depth, but
   of wavy character, in which the _sides_ of the groove act upon the
   tracing point instead of the bottom, as is the case in the vertical
   type. This form we named the zig-zag form, and referred to it in
   that way in our notes. Its important advantage in guiding the
   reproducing needle I first called attention to in the note on p.
   9-Vol 1-Home Notes on March 29-1881, and endeavored to use it in my
   early work, but encountered so much difficulty in getting a form of
   reproducer that would work with the soft wax records without
   tearing the groove, we used the hill and valley type of record more
   often than the other.
  \end{verbatim}
  
  In 1885, when the Volta associates were sure that they had a number of
  practical inventions, they filed applications for patents. They also
  began to look around for investors. After giving several demonstrations
  in Washington, they gained the necessary support, and the American
  Graphophone Co. was organized to manufacture and sell the machines. The
  Volta Graphophone Co. was formed to control the patents.
  
  The Howe sewing machine factory at Bridgeport, Connecticut, became the
  American Graphophone plant; Tainter went there to supervise the
  manufacturing, and continued his inventive work for many years. This
  Bridgeport plant is still in use today by a successor firm, the
  Dictaphone Corporation.
  
  The work of the Volta associates laid the foundation for the successful
  use of the dictating machine in business, for their wax recording
  process was practical and their machines sturdy. But it was to take
  several more years and the renewed work of Edison and further
  developments by Berliner and many others, before the talking machine
  industry really got under way and became a major factor in home
  entertainment.{[}10{]}
  
  PATENTS WHICH RESULTED FROM THE VOLTA LABORATORY ASSOCIATION
  
  \emph{Patent Number} \emph{Year} \emph{Patent} \emph{Inventors} 229495
  1880 Telephone call register C. S. Tainter 235496 1880 Photophone
  transmitter A. G. Bell, C. S. Tainter 235497 1880 Selenium cells A. G.
  Bell, C. S. Tainter 235590 1800 Selenium cells C. S. Tainter 241909 1881
  Photophonic receiver A. G. Bell, C. S. Tainter 243657 1881 Telephone
  transmitter C. S. Tainter 289725 1883 Electric conductor C. S. Tainter
  336081 1886 Transmitter for electric C. A. Bell telephone lines 336082
  1886 Jet microphone for C. A. Bell transmitting sounds by means of jets
  336083 1886 Telephone transmitter C. A. Bell 336173 1886 Telephone
  transmitter C. S. Tainter 341212 1886 Reproducing sounds A. G. Bell,
  from phonograph C. A. Bell, records C. S. Tainter 341213 1886
  Reproducing and A. G. Bell, recording sounds by C. A. Bell, radiant
  energy C. S. Tainter 341214 1886 Recording and reproducing C. A. Bell,
  speech and other sounds C. S. Tainter 341287 1886 Recording and
  reproducing C. S. Tainter sounds 341288 1886 Apparatus for recording C.
  S. Tainter and reproducing sounds 374133 1887 Paper cylinder for C. S.
  Tainter graphophonic records 375579 1887 Apparatus for recording C. S.
  Tainter and reproducing speech and other sounds 380535 1888 Graphophone
  C. S. Tainter 385886 1888 Graphophone C. S. Tainter 385887 1888
  Graphophonic tablet C. S. Tainter 388462 1888 Machine for making C. S.
  Tainter paper tubes 392763 1888 Mounting for diaphragms C. S. Tainter
  for acoustical instruments 393190 1888 Tablet for use in C. S. Tainter
  graphophones 393191 1888 Support for graphophonic C. S. Tainter tablets
  416969 1889 Speed regulator C. S. Tainter 421450 1890 Graphophone tablet
  C. S. Tainter 428646 1890 Machine for the manufacture C. S. Tainter of
  wax coated tablets for graphophones 506348 1893 Coin controlled
  graphophone C. S. Tainter 510656 1893 Reproducer for graphophones C. S.
  Tainter 670442 1901 Graphophone record C. S. Tainter duplicating machine
  730986 1903 Graphophone C. S. Tainter
  
  CONTENTS OF SMALL CHEST RECEIVED BY THE SMITHSONIAN INSTITUTION FROM
  MRS. LAURA F. TAINTER, 1947
  
  Books (10) containing the home notes, volumes 1 to 8 and 11 and 12,
  inclusive, March 1881-November 1883. (Vols, 9, 10, and 13 were burned in
  a laboratory fire, September 1897.)
  
  Binder containing drawings and notes for multiple record duplicator,
  October 8, 1897-1908, and miscellaneous inquiries, log, telegraph
  recorder, diet, home plans.
  
  Binder containing printed specifications of patents, S. Tainter, A. G.
  Bell, and C. A. Bell, June 29, 1880 to June 16, 1903.
  
  Medal, Exposition Internationale d'Electricite, Paris, 1881, marked
  ``Tainter.''
  
  Medal, Panama-Pacific Exposition, San Francisco, 1915, Medal of Award.
  
  Seven purple lapel rosettes (?), one with ribbon and palms, in boxes
  marked ``1890.'' Notes in newspaper clipping.
  
  Records of testimony of C. S. Tainter in various suits involving the
  phonograph: Volta Graphophone Co. \emph{vs.} Columbia Phonograph Co.,
  no. 14533, Supreme Court of the District of Columbia, dated January 13,
  1894; American Graphophone Co. \emph{vs.} U. S. Phonograph Co., U. S.
  Circuit Court, New Jersey, dated May 14, 1895; American Graphophone Co.
  \emph{vs.} Edw. H. Amet, U. S. Circuit Court, Northern District,
  Illinois, in equity, dated February 14, 1896; American Graphophone Co.
  \emph{vs.} U. S. Phonograph Co., \emph{et al.}, U. S. Circuit Court, New
  Jersey, in equity, no date; American Graphophone Co. \emph{vs.} Leeds
  \emph{et al.}, U. S. Circuit Court, Two District, New York, N. Y., no
  date; testimony marked ``Questions asked in Edison Co. suits''
  (duplicate copies) no date, no citation.
  
  CONTENTS OF SMALL CHEST RECEIVED BY THE SMITHSONIAN INSTITUTION FROM
  MRS. LAURA F. TAINTER, 1950
  
  Typed manuscript--``Memoirs of Charles Sumner Tainter'' (plus many
  photostats of notes and articles) 4-1/2 inches thick, pp.~1-71 to about
  1878, pp.~1 to 104 to factory at Bridgeport, some pages missing.
  
  Box--containing handwritten notes for ``memoirs'' includes copies of
  text of above (less photostats); copies of short biography; agreement
  creating American Graphophone Co.; letter of election to life membership
  in the American Association for the Advancement of Science.
  
  Binder--exhibits of Tainter drawings in American Graphophone Co.
  \emph{vs.} Edison Phonograph Works., vol.~1, U. S. Circuit Court, New
  Jersey.
  
  Folder--clippings and photostats relating to the machines deposited in
  Smithsonian.
  
  Certificate of appointment ``Officer de l'Instruction publique,''
  France, October 31, 1889, for exhibition of Graphophone, Exposition
  Universelle, 1889.
  
  Framed photo of Berliner \& Tainter, 1919.
  
  Photo of Tainter, 1919.
  
  Separate package containing gold medal, certificate, Panama-Pacific
  Exposition, San Francisco, 1915; gold medal, certificate, Exposition
  Internationale Electricite, Paris, 1881.
  
  \section*{Footnotes}\label{footnotes}
  \addcontentsline{toc}{section}{Footnotes}
  
  {[}1{]} Charles Sumner Tainter (1854-1940), ``The talking machine and
  some little known facts in connection with its early development,''
  unpublished manuscript in the collections of the U. S. National Museum.
  
  {[}2{]} One of the most interesting prophecies was written in 1656 by
  Cyrano de Bergerac, in his \emph{Comic history of the states and empires
  of the Moon}:
  
  ``'I began to study closely my books and their covers which impressed me
  for their richness. One was decorated with a single diamond, more
  brilliant by far than ours. The second seemed but a single pearl cleft
  in twain.
  
  ``\,`When I opened the covers, I found inside something made of metal,
  not unlike our clocks, full of mysterious little springs and almost
  invisible mechanisms. 'Tis a book, 'tis true, but a miraculous book,
  which has no pages or letters. Indeed, 'tis a book which to enjoy the
  eyes are useless; only ears suffice. When a man desires to read, then,
  he surrounds this contrivance with many small tendons of every kind,
  then he places the needle on the chapter to be heard and, at the same
  time, there come, as from the mouth of a man or from an instrument of
  music, all those clear and separate sounds which make up the Lunarians'
  tongue.''' (See A. Coeuroy and G. Clarence, \emph{Le phonographe},
  Paris, 1929, p.~9, 10.)
  
  {[}3{]} Tainter retained a lifelong admiration for Alexander Graham
  Bell. This is Tainter's description of their first meeting in Cambridge:
  ``\ldots{} one day I received a visit from a very distinguished looking
  gentleman with jet black hair and beard, who announced himself as Mr.~A.
  Graham Bell. His charm of manner and conversation attracted me
  greatly\ldots{}.'' Tainter, \emph{op. cit.} (footnote 1), p.~2.
  
  {[}4{]} A. G. Bell apparently spent little time in the Volta Laboratory.
  The Dr.~Bell referred to in Tainter's notebooks is Chichester A. Bell.
  The basic graphophone patent (U. S. patent 341214) was issued to C. A.
  Bell and Tainter. The Tainter material reveals A. G. Bell as the man who
  suggested the basic lines of research (and furnished the money), and
  then allowed his associates to get the credit for many of the inventions
  that resulted.
  
  {[}5{]} Tainter, \emph{op. cit.} (footnote 1), p.~3.
  
  {[}6{]} \_Ibid.\_, p.~5.
  
  {[}7{]} \_Ibid.\_, p.~30.
  
  {[}8{]} As quoted by \emph{The Washington Herald}, October 28, 1937.
  
  {[}9{]} Tainter, \emph{op. cit.} (footnote 1), pp.~28, 29.
  
  {[}10{]} The basic distinction between the first Edison patent, and the
  Bell and Tainter patent of 1886 was the method of recording. Edison's
  method was to \emph{indent} the sound waves on a piece of tin-foil (wax
  was included as a recording material in his English patent); the Bell
  and Tainter improvement called for \emph{cutting} or
  ``\emph{engraving}'' the sound waves into a wax record, with a sharp
  recording stylus.
  
  The strength of Bell and Tainter patent is indicated by the following
  excerpt from a letter written by a Washington patent attorney, S. T.
  Cameron, who was a member of the law firm which carried on litigation
  for the American Graphophone Co. The letter is dated December 8, 1914,
  and is addressed to George C. Maynard, Curator of Mechanical Technology,
  U. S. National Museum: ``Subsequent to the issuance of the Bell and
  Tainter patent No. 341214, Edison announced that he would shortly
  produce his `new phonograph' which, when it appeared, was in fact
  nothing but the Bell and Tainter record set forth in their patent
  341214, being a record cut or engraved in wax or wax-like material,
  although Edison always insisted on calling this record an `indented'
  record, doubtless because his original tin-foil record was an `indented'
  record. Edison was compelled to acknowledge that his `new phonograph'
  was an infringement of the Bell and Tainter patent 341214, and took out
  a license under the Bell and Tainter patent and made his records under
  that patent as the result of that license.''
  
  \section*{Postscript}\label{postscript}
  \addcontentsline{toc}{section}{Postscript}
  
  \textbf{\emph{END OF THE PROJECT GUTENBERG EBOOK DEVELOPMENT OF THE
  PHONOGRAPH AT ALEXANDER GRAHAM BELL'S VOLTA LABORATORY}}
  
  ******* This file should be named 30112.txt or 30112.zip *******
  
  This and all associated files of various formats will be found in:
  http://www.gutenberg.org/dirs/3/0/1/1/30112
  
  Updated editions will replace the previous one--the old editions will be
  renamed.
  
  Creating the works from public domain print editions means that no one
  owns a United States copyright in these works, so the Foundation (and
  you!) can copy and distribute it in the United States without permission
  and without paying copyright royalties. Special rules, set forth in the
  General Terms of Use part of this license, apply to copying and
  distributing Project Gutenberg-tm electronic works to protect the
  PROJECT GUTENBERG-tm concept and trademark. Project Gutenberg is a
  registered trademark, and may not be used if you charge for the eBooks,
  unless you receive specific permission. If you do not charge anything
  for copies of this eBook, complying with the rules is very easy. You may
  use this eBook for nearly any purpose such as creation of derivative
  works, reports, performances and research. They may be modified and
  printed and given away--you may do practically ANYTHING with public
  domain eBooks. Redistribution is subject to the trademark license,
  especially commercial redistribution.
  
  *** START: FULL LICENSE ***
  
  THE FULL PROJECT GUTENBERG LICENSE PLEASE READ THIS BEFORE YOU
  DISTRIBUTE OR USE THIS WORK
  
  To protect the Project Gutenberg-tm mission of promoting the free
  distribution of electronic works, by using or distributing this work (or
  any other work associated in any way with the phrase ``Project
  Gutenberg''), you agree to comply with all the terms of the Full Project
  Gutenberg-tm License (available with this file or online at
  http://www.gutenberg.org/license).
  
  Section 1. General Terms of Use and Redistributing Project Gutenberg-tm
  electronic works
  
  1.A. By reading or using any part of this Project Gutenberg-tm
  electronic work, you indicate that you have read, understand, agree to
  and accept all the terms of this license and intellectual property
  (trademark/copyright) agreement. If you do not agree to abide by all the
  terms of this agreement, you must cease using and return or destroy all
  copies of Project Gutenberg-tm electronic works in your possession. If
  you paid a fee for obtaining a copy of or access to a Project
  Gutenberg-tm electronic work and you do not agree to be bound by the
  terms of this agreement, you may obtain a refund from the person or
  entity to whom you paid the fee as set forth in paragraph 1.E.8.
  
  1.B. ``Project Gutenberg'' is a registered trademark. It may only be
  used on or associated in any way with an electronic work by people who
  agree to be bound by the terms of this agreement. There are a few things
  that you can do with most Project Gutenberg-tm electronic works even
  without complying with the full terms of this agreement. See paragraph
  1.C below. There are a lot of things you can do with Project
  Gutenberg-tm electronic works if you follow the terms of this agreement
  and help preserve free future access to Project Gutenberg-tm electronic
  works. See paragraph 1.E below.
  
  1.C. The Project Gutenberg Literary Archive Foundation (``the
  Foundation'' or PGLAF), owns a compilation copyright in the collection
  of Project Gutenberg-tm electronic works. Nearly all the individual
  works in the collection are in the public domain in the United States.
  If an individual work is in the public domain in the United States and
  you are located in the United States, we do not claim a right to prevent
  you from copying, distributing, performing, displaying or creating
  derivative works based on the work as long as all references to Project
  Gutenberg are removed. Of course, we hope that you will support the
  Project Gutenberg-tm mission of promoting free access to electronic
  works by freely sharing Project Gutenberg-tm works in compliance with
  the terms of this agreement for keeping the Project Gutenberg-tm name
  associated with the work. You can easily comply with the terms of this
  agreement by keeping this work in the same format with its attached full
  Project Gutenberg-tm License when you share it without charge with
  others.
  
  1.D. The copyright laws of the place where you are located also govern
  what you can do with this work. Copyright laws in most countries are in
  a constant state of change. If you are outside the United States, check
  the laws of your country in addition to the terms of this agreement
  before downloading, copying, displaying, performing, distributing or
  creating derivative works based on this work or any other Project
  Gutenberg-tm work. The Foundation makes no representations concerning
  the copyright status of any work in any country outside the United
  States.
  
  1.E. Unless you have removed all references to Project Gutenberg:
  
  1.E.1. The following sentence, with active links to, or other immediate
  access to, the full Project Gutenberg-tm License must appear prominently
  whenever any copy of a Project Gutenberg-tm work (any work on which the
  phrase ``Project Gutenberg'' appears, or with which the phrase ``Project
  Gutenberg'' is associated) is accessed, displayed, performed, viewed,
  copied or distributed:
  
  This eBook is for the use of anyone anywhere at no cost and with almost
  no restrictions whatsoever. You may copy it, give it away or re-use it
  under the terms of the Project Gutenberg License included with this
  eBook or online at www.gutenberg.org
  
  1.E.2. If an individual Project Gutenberg-tm electronic work is derived
  from the public domain (does not contain a notice indicating that it is
  posted with permission of the copyright holder), the work can be copied
  and distributed to anyone in the United States without paying any fees
  or charges. If you are redistributing or providing access to a work with
  the phrase ``Project Gutenberg'' associated with or appearing on the
  work, you must comply either with the requirements of paragraphs 1.E.1
  through 1.E.7 or obtain permission for the use of the work and the
  Project Gutenberg-tm trademark as set forth in paragraphs 1.E.8 or
  1.E.9.
  
  1.E.3. If an individual Project Gutenberg-tm electronic work is posted
  with the permission of the copyright holder, your use and distribution
  must comply with both paragraphs 1.E.1 through 1.E.7 and any additional
  terms imposed by the copyright holder. Additional terms will be linked
  to the Project Gutenberg-tm License for all works posted with the
  permission of the copyright holder found at the beginning of this work.
  
  1.E.4. Do not unlink or detach or remove the full Project Gutenberg-tm
  License terms from this work, or any files containing a part of this
  work or any other work associated with Project Gutenberg-tm.
  
  1.E.5. Do not copy, display, perform, distribute or redistribute this
  electronic work, or any part of this electronic work, without
  prominently displaying the sentence set forth in paragraph 1.E.1 with
  active links or immediate access to the full terms of the Project
  Gutenberg-tm License.
  
  1.E.6. You may convert to and distribute this work in any binary,
  compressed, marked up, nonproprietary or proprietary form, including any
  word processing or hypertext form. However, if you provide access to or
  distribute copies of a Project Gutenberg-tm work in a format other than
  ``Plain Vanilla ASCII'' or other format used in the official version
  posted on the official Project Gutenberg-tm web site
  (www.gutenberg.org), you must, at no additional cost, fee or expense to
  the user, provide a copy, a means of exporting a copy, or a means of
  obtaining a copy upon request, of the work in its original ``Plain
  Vanilla ASCII'' or other form. Any alternate format must include the
  full Project Gutenberg-tm License as specified in paragraph 1.E.1.
  
  1.E.7. Do not charge a fee for access to, viewing, displaying,
  performing, copying or distributing any Project Gutenberg-tm works
  unless you comply with paragraph 1.E.8 or 1.E.9.
  
  1.E.8. You may charge a reasonable fee for copies of or providing access
  to or distributing Project Gutenberg-tm electronic works provided that
  
  \begin{itemize}
  \item
    You pay a royalty fee of 20\% of the gross profits you derive from the
    use of Project Gutenberg-tm works calculated using the method you
    already use to calculate your applicable taxes. The fee is owed to the
    owner of the Project Gutenberg-tm trademark, but he has agreed to
    donate royalties under this paragraph to the Project Gutenberg
    Literary Archive Foundation. Royalty payments must be paid within 60
    days following each date on which you prepare (or are legally required
    to prepare) your periodic tax returns. Royalty payments should be
    clearly marked as such and sent to the Project Gutenberg Literary
    Archive Foundation at the address specified in Section 4,
    ``Information about donations to the Project Gutenberg Literary
    Archive Foundation.''
  \item
    You provide a full refund of any money paid by a user who notifies you
    in writing (or by e-mail) within 30 days of receipt that s/he does not
    agree to the terms of the full Project Gutenberg-tm License. You must
    require such a user to return or destroy all copies of the works
    possessed in a physical medium and discontinue all use of and all
    access to other copies of Project Gutenberg-tm works.
  \item
    You provide, in accordance with paragraph 1.F.3, a full refund of any
    money paid for a work or a replacement copy, if a defect in the
    electronic work is discovered and reported to you within 90 days of
    receipt of the work.
  \item
    You comply with all other terms of this agreement for free
    distribution of Project Gutenberg-tm works.
  \end{itemize}
  
  1.E.9. If you wish to charge a fee or distribute a Project Gutenberg-tm
  electronic work or group of works on different terms than are set forth
  in this agreement, you must obtain permission in writing from both the
  Project Gutenberg Literary Archive Foundation and Michael Hart, the
  owner of the Project Gutenberg-tm trademark. Contact the Foundation as
  set forth in Section 3 below.
  
  1.F.
  
  1.F.1. Project Gutenberg volunteers and employees expend considerable
  effort to identify, do copyright research on, transcribe and proofread
  public domain works in creating the Project Gutenberg-tm collection.
  Despite these efforts, Project Gutenberg-tm electronic works, and the
  medium on which they may be stored, may contain ``Defects,'' such as,
  but not limited to, incomplete, inaccurate or corrupt data,
  transcription errors, a copyright or other intellectual property
  infringement, a defective or damaged disk or other medium, a computer
  virus, or computer codes that damage or cannot be read by your
  equipment.
  
  1.F.2. LIMITED WARRANTY, DISCLAIMER OF DAMAGES - Except for the ``Right
  of Replacement or Refund'' described in paragraph 1.F.3, the Project
  Gutenberg Literary Archive Foundation, the owner of the Project
  Gutenberg-tm trademark, and any other party distributing a Project
  Gutenberg-tm electronic work under this agreement, disclaim all
  liability to you for damages, costs and expenses, including legal fees.
  YOU AGREE THAT YOU HAVE NO REMEDIES FOR NEGLIGENCE, STRICT LIABILITY,
  BREACH OF WARRANTY OR BREACH OF CONTRACT EXCEPT THOSE PROVIDED IN
  PARAGRAPH F3. YOU AGREE THAT THE FOUNDATION, THE TRADEMARK OWNER, AND
  ANY DISTRIBUTOR UNDER THIS AGREEMENT WILL NOT BE LIABLE TO YOU FOR
  ACTUAL, DIRECT, INDIRECT, CONSEQUENTIAL, PUNITIVE OR INCIDENTAL DAMAGES
  EVEN IF YOU GIVE NOTICE OF THE POSSIBILITY OF SUCH DAMAGE.
  
  1.F.3. LIMITED RIGHT OF REPLACEMENT OR REFUND - If you discover a defect
  in this electronic work within 90 days of receiving it, you can receive
  a refund of the money (if any) you paid for it by sending a written
  explanation to the person you received the work from. If you received
  the work on a physical medium, you must return the medium with your
  written explanation. The person or entity that provided you with the
  defective work may elect to provide a replacement copy in lieu of a
  refund. If you received the work electronically, the person or entity
  providing it to you may choose to give you a second opportunity to
  receive the work electronically in lieu of a refund. If the second copy
  is also defective, you may demand a refund in writing without further
  opportunities to fix the problem.
  
  1.F.4. Except for the limited right of replacement or refund set forth
  in paragraph 1.F.3, this work is provided to you `AS-IS', WITH NO OTHER
  WARRANTIES OF ANY KIND, EXPRESS OR IMPLIED, INCLUDING BUT NOT LIMITED TO
  WARRANTIES OF MERCHANTIBILITY OR FITNESS FOR ANY PURPOSE.
  
  1.F.5. Some states do not allow disclaimers of certain implied
  warranties or the exclusion or limitation of certain types of damages.
  If any disclaimer or limitation set forth in this agreement violates the
  law of the state applicable to this agreement, the agreement shall be
  interpreted to make the maximum disclaimer or limitation permitted by
  the applicable state law. The invalidity or unenforceability of any
  provision of this agreement shall not void the remaining provisions.
  
  1.F.6. INDEMNITY - You agree to indemnify and hold the Foundation, the
  trademark owner, any agent or employee of the Foundation, anyone
  providing copies of Project Gutenberg-tm electronic works in accordance
  with this agreement, and any volunteers associated with the production,
  promotion and distribution of Project Gutenberg-tm electronic works,
  harmless from all liability, costs and expenses, including legal fees,
  that arise directly or indirectly from any of the following which you do
  or cause to occur: (a) distribution of this or any Project Gutenberg-tm
  work, (b) alteration, modification, or additions or deletions to any
  Project Gutenberg-tm work, and (c) any Defect you cause.
  
  Section 2. Information about the Mission of Project Gutenberg-tm
  
  Project Gutenberg-tm is synonymous with the free distribution of
  electronic works in formats readable by the widest variety of computers
  including obsolete, old, middle-aged and new computers. It exists
  because of the efforts of hundreds of volunteers and donations from
  people in all walks of life.
  
  Volunteers and financial support to provide volunteers with the
  assistance they need are critical to reaching Project Gutenberg-tm's
  goals and ensuring that the Project Gutenberg-tm collection will remain
  freely available for generations to come. In 2001, the Project Gutenberg
  Literary Archive Foundation was created to provide a secure and
  permanent future for Project Gutenberg-tm and future generations. To
  learn more about the Project Gutenberg Literary Archive Foundation and
  how your efforts and donations can help, see Sections 3 and 4 and the
  Foundation web page at http://www.gutenberg.org/fundraising/pglaf.
  
  Section 3. Information about the Project Gutenberg Literary Archive
  Foundation
  
  The Project Gutenberg Literary Archive Foundation is a non profit
  501(c)(3) educational corporation organized under the laws of the state
  of Mississippi and granted tax exempt status by the Internal Revenue
  Service. The Foundation's EIN or federal tax identification number is
  64-6221541. Contributions to the Project Gutenberg Literary Archive
  Foundation are tax deductible to the full extent permitted by U.S.
  federal laws and your state's laws.
  
  The Foundation's principal office is located at 4557 Melan Dr.~S.
  Fairbanks, AK, 99712., but its volunteers and employees are scattered
  throughout numerous locations. Its business office is located at 809
  North 1500 West, Salt Lake City, UT 84116, (801) 596-1887, email
  business@pglaf.org. Email contact links and up to date contact
  information can be found at the Foundation's web site and official page
  at http://www.gutenberg.org/about/contact
  
  For additional contact information: Dr.~Gregory B. Newby Chief Executive
  and Director gbnewby@pglaf.org
  
  Section 4. Information about Donations to the Project Gutenberg Literary
  Archive Foundation
  
  Project Gutenberg-tm depends upon and cannot survive without wide spread
  public support and donations to carry out its mission of increasing the
  number of public domain and licensed works that can be freely
  distributed in machine readable form accessible by the widest array of
  equipment including outdated equipment. Many small donations (\$1 to
  \$5,000) are particularly important to maintaining tax exempt status
  with the IRS.
  
  The Foundation is committed to complying with the laws regulating
  charities and charitable donations in all 50 states of the United
  States. Compliance requirements are not uniform and it takes a
  considerable effort, much paperwork and many fees to meet and keep up
  with these requirements. We do not solicit donations in locations where
  we have not received written confirmation of compliance. To SEND
  DONATIONS or determine the status of compliance for any particular state
  visit http://www.gutenberg.org/fundraising/donate
  
  While we cannot and do not solicit contributions from states where we
  have not met the solicitation requirements, we know of no prohibition
  against accepting unsolicited donations from donors in such states who
  approach us with offers to donate.
  
  International donations are gratefully accepted, but we cannot make any
  statements concerning tax treatment of donations received from outside
  the United States. U.S. laws alone swamp our small staff.
  
  Please check the Project Gutenberg Web pages for current donation
  methods and addresses. Donations are accepted in a number of other ways
  including checks, online payments and credit card donations. To donate,
  please visit: http://www.gutenberg.org/fundraising/donate
  
  Section 5. General Information About Project Gutenberg-tm electronic
  works.
  
  Professor Michael S. Hart is the originator of the Project Gutenberg-tm
  concept of a library of electronic works that could be freely shared
  with anyone. For thirty years, he produced and distributed Project
  Gutenberg-tm eBooks with only a loose network of volunteer support.
  
  Project Gutenberg-tm eBooks are often created from several printed
  editions, all of which are confirmed as Public Domain in the U.S. unless
  a copyright notice is included. Thus, we do not necessarily keep eBooks
  in compliance with any particular paper edition.
  
  Most people start at our Web site which has the main PG search facility:
  
  \begin{verbatim}
   http://www.gutenberg.org
  \end{verbatim}
  
  This Web site includes information about Project Gutenberg-tm, including
  how to make donations to the Project Gutenberg Literary Archive
  Foundation, how to help produce our new eBooks, and how to subscribe to
  our email newsletter to hear about new eBooks.
  
\end{document}
